% Options for packages loaded elsewhere
\PassOptionsToPackage{unicode}{hyperref}
\PassOptionsToPackage{hyphens}{url}
\documentclass[
]{article}
\usepackage{xcolor}
\usepackage{amsmath,amssymb}
\setcounter{secnumdepth}{-\maxdimen} % remove section numbering
\usepackage{iftex}
\ifPDFTeX
  \usepackage[T1]{fontenc}
  \usepackage[utf8]{inputenc}
  \usepackage{textcomp} % provide euro and other symbols
\else % if luatex or xetex
  \usepackage{unicode-math} % this also loads fontspec
  \defaultfontfeatures{Scale=MatchLowercase}
  \defaultfontfeatures[\rmfamily]{Ligatures=TeX,Scale=1}
\fi
\usepackage{lmodern}
\ifPDFTeX\else
  % xetex/luatex font selection
\fi
% Use upquote if available, for straight quotes in verbatim environments
\IfFileExists{upquote.sty}{\usepackage{upquote}}{}
\IfFileExists{microtype.sty}{% use microtype if available
  \usepackage[]{microtype}
  \UseMicrotypeSet[protrusion]{basicmath} % disable protrusion for tt fonts
}{}
\makeatletter
\@ifundefined{KOMAClassName}{% if non-KOMA class
  \IfFileExists{parskip.sty}{%
    \usepackage{parskip}
  }{% else
    \setlength{\parindent}{0pt}
    \setlength{\parskip}{6pt plus 2pt minus 1pt}}
}{% if KOMA class
  \KOMAoptions{parskip=half}}
\makeatother
\usepackage{longtable,booktabs,array}
\usepackage{caption}
\captionsetup[table]{skip=6pt}
\newcounter{none} % for unnumbered tables
\usepackage{calc} % for calculating minipage widths
% Correct order of tables after \paragraph or \subparagraph
\usepackage{etoolbox}
\makeatletter
\patchcmd\longtable{\par}{\if@noskipsec\mbox{}\fi\par}{}{}
\makeatother
% Allow footnotes in longtable head/foot
\IfFileExists{footnotehyper.sty}{\usepackage{footnotehyper}}{\usepackage{footnote}}
\makesavenoteenv{longtable}
\usepackage{graphicx}
\makeatletter
\newsavebox\pandoc@box
\newcommand*\pandocbounded[1]{% scales image to fit in text height/width
  \sbox\pandoc@box{#1}%
  \Gscale@div\@tempa{\textheight}{\dimexpr\ht\pandoc@box+\dp\pandoc@box\relax}%
  \Gscale@div\@tempb{\linewidth}{\wd\pandoc@box}%
  \ifdim\@tempb\p@<\@tempa\p@\let\@tempa\@tempb\fi% select the smaller of both
  \ifdim\@tempa\p@<\p@\scalebox{\@tempa}{\usebox\pandoc@box}%
  \else\usebox{\pandoc@box}%
  \fi%
}
% Set default figure placement to htbp
\def\fps@figure{htbp}
\makeatother
\setlength{\emergencystretch}{3em} % prevent overfull lines
\providecommand{\tightlist}{%
  \setlength{\itemsep}{0pt}\setlength{\parskip}{0pt}}
\usepackage{bookmark}
\IfFileExists{xurl.sty}{\usepackage{xurl}}{} % add URL line breaks if available
\urlstyle{same}
% fallback for those not using the hyperref driver hyperxmp:
\makeatletter
\@ifundefined{xmpquote}{\newcommand{\xmpquote}[1]{#1}}{}
\makeatother
\hypersetup{
  hidelinks,
  pdfcreator={LaTeX via pandoc}}

\author{}
\date{}

\begin{document}

\section{Kualitas Lingkungan di Kawasan
Smelter}\label{kualitas-lingkungan-di-kawasan-smelter}

\emph{Menguji secara empiris korelasi antara intensitas ekspansi
fasilitas peleburan nikel (smelter) dengan anjloknya Indeks Kualitas Air
(IKA), Udara (IKU), dan laju deforestasi komoditas di Pulau Sulawesi.}

\begin{quote}
\textbf{Alur Kausalitas (Ekonomi Politik Ekologi):}
\texttt{Konsentrasi\ Smelter\ \&\ PLTU} →
\texttt{Pembuangan\ Tailing\ \&\ Emisi\ Partikulat} →
\texttt{Pencemaran\ Air\ \&\ Udara} →
\texttt{Daya\ Dukung\ Lingkungan\ Kritis}

Industrialisasi ekstraktif bekerja seperti parasit terhadap daya tampung
lingkungan. Fasilitas peleburan membutuhkan daya luar biasa besar
(dipasok PLTU batu bara) dan menghasilkan limbah B3 (slag/tailing) dalam
jumlah raksasa. Kehadiran industri ini secara absolut menurunkan ambang
batas toleransi alam, yang direkam melalui indikator baku mutu nasional.

\textbf{Variabel Tekanan (X):} * \textbf{Jumlah Smelter \& PLTU
Captive:} Konsentrasi fasilitas peleburan dan pembangkit batu bara
(ESDM, GEM). * \textbf{Luas Kawasan Industri:} Ekspansi spasial proyek
strategis.

\textbf{Variabel Dampak Ekologis (Y):} * \textbf{Indeks Kualitas Air
(IKA):} Skor kualitas air berdasarkan parameter fisik/kimia (KLHK, BPS).
* \textbf{Indeks Kualitas Udara (IKU):} Skor pencemaran udara ambien
(KLHK, BPS). * \textbf{Laju Deforestasi Komoditas:} Kehilangan tutupan
pohon akibat kegiatan ekstraktif (Global Forest Watch).

\textbf{Metode Pengolahan Data:} Analisis menggunakan pendekatan
\emph{Cross-sectional} dan \emph{Time-Series} (Panel Data). Korelasi
dibuktikan secara statistik melalui uji \textbf{Crosstabulation
(Chi-Square/Symmetric Measures)} untuk melihat seberapa jauh status
`Kritis' pada wilayah beririsan dengan label `Tinggi' pada kehadiran
industri.
\end{quote}

\subsection{Tumbal Ekologis: Ketika 778 Smelter Mencekik Napas dan Air
Sulawesi}\label{tumbal-ekologis-ketika-778-smelter-mencekik-napas-dan-air-sulawesi}

Pesta ``Hilirisasi'' nikel ternyata menagih bayaran tunai langsung ke
jantung ekosistem Pulau Sulawesi. Konsentrasi \textbf{778 fasilitas
mega-smelter} yang ditenagai oleh mesin pembakaran batu bara raksasa
sebesar \textbf{9,825 MW} tidak hanya mencaplok ruang spasial, tetapi
telah membengkokkan kurva daya dukung lingkungan ke titik terendahnya.
Di bawah bayang-bayang narasi transisi energi, kualitas hidup jutaan
warga dipertaruhkan oleh penurunan masif kualitas air dan udara,
ditambah ancaman ledakan limbah beracun.

Data empiris merekam kebangkrutan ini secara presisi. Operasional
industri ekstraktif telah merobek \textbf{2,107,041 Hektar} ruang hidup
alami dan menumpuk lebih dari \textbf{33.8 Juta Ton} limbah B3/tailing
per tahun. Beban polusi ini membuat Indeks Kualitas Air (IKA) rata-rata
merosot ke angka \textbf{58.8}, dan Indeks Kualitas Udara (IKU) anjlok
ke level \textbf{92.4}. Angka ini lebih dari sekadar indikator; ia
adalah bukti forensik kegagalan kebijakan perlindungan ekologis negara
di hadapan invasi modal korporasi.

\subsubsection{Metrik Kritis Kualitas
Lingkungan}\label{metrik-kritis-kualitas-lingkungan}

{\def\LTcaptype{none} % do not increment counter
\begin{longtable}[]{@{}
  >{\raggedright\arraybackslash}p{(\linewidth - 6\tabcolsep) * \real{0.2500}}
  >{\raggedright\arraybackslash}p{(\linewidth - 6\tabcolsep) * \real{0.2500}}
  >{\raggedright\arraybackslash}p{(\linewidth - 6\tabcolsep) * \real{0.2500}}
  >{\raggedright\arraybackslash}p{(\linewidth - 6\tabcolsep) * \real{0.2500}}@{}}
\toprule\noalign{}
\begin{minipage}[b]{\linewidth}\raggedright
Indikator
\end{minipage} & \begin{minipage}[b]{\linewidth}\raggedright
Nilai
\end{minipage} & \begin{minipage}[b]{\linewidth}\raggedright
Deskripsi
\end{minipage} & \begin{minipage}[b]{\linewidth}\raggedright
Sumber
\end{minipage} \\
\midrule\noalign{}
\endhead
\bottomrule\noalign{}
\endlastfoot
\textbf{Indeks Kualitas Air (2023)} & \textbf{58.8 Poin} & Tren
penurunan IKA di Sulawesi akibat cemaran logam berat dan sedimen operasi
tambang. & KLHK \& BPS (SLHI) \\
\textbf{Polusi Udara NO₂ (NASA)} & \textbf{5.76e-06 mol/m²} & Satelit
TROPOMI membongkar paradoks IKU resmi. Konsentrasi gas beracun meroket
tajam seiring ekspansi PLTU captive. & Satelit Sentinel-5P (Google Earth
Engine) \\
\textbf{Timbulan Limbah B3} & \textbf{33.8 Jt Ton} & Estimasi produksi
limbah tailing dan slag per tahun dari kawasan mega-industri di
Sulawesi. & Data Ekstraksi NGO \& AMDAL \\
\textbf{Konversi Deforestasi} & \textbf{2,107,041 Ha} & Luasan tutupan
hutan yang hancur dibabat untuk pembukaan lubang tambang nikel. & Global
Forest Watch (GFW) \\
\end{longtable}
}

\begin{center}\rule{0.5\linewidth}{0.5pt}\end{center}

\subsubsection{2.1. Dampak Limbah Tailing: Konsentrasi Smelter vs Indeks
Kualitas Air
(IKA)}\label{dampak-limbah-tailing-konsentrasi-smelter-vs-indeks-kualitas-air-ika}

\begin{quote}
\textbf{Metode Analisis:} Sub-bab ini menggunakan pendekatan Analisis
Spasial dan Uji Statistik Chi-Square (Crosstabulation) untuk mengukur
dampak konsentrasi smelter terhadap penurunan kualitas air.

\begin{enumerate}
\def\labelenumi{\arabic{enumi}.}
\tightlist
\item
  \textbf{Uji Tabulasi Silang (Chi-Square Test of Independence):}
  Binning Kategori via Median. Decision Rule: P-Value \textless{} 0.05
  maka Tolak H0.
\item
  \textbf{Variabel \& Fitur Data:} Jumlah\_Smelter (X), Indeks Kualitas
  Air (Y), Data Panel 2016-2023.
\end{enumerate}
\end{quote}

Tingginya aktivitas pengolahan bijih nikel (\emph{smelter}) secara
langsung memicu ekskresi limbah \emph{tailing} dan terak (\emph{slag})
dalam skala masif. Peta geospasial dan agregasi data di bawah ini tidak
sekadar menunjukkan visualisasi statis, melainkan menyingkap realitas
berdarah dari proyek ambisius hilirisasi nikel yang dipaksakan di
Sulawesi. Dari total \textbf{778 fasilitas mega-smelter} yang
beroperasi, konsentrasi absolut pembongkaran ruang terpusat pada dua
episentrum industri, yakni Sulawesi Tengah dengan \textbf{344 fasilitas
smelter} dan Sulawesi Tenggara dengan \textbf{262 fasilitas smelter}.
Ekspansi agresif yang tidak terkendali ini memonopoli pemanfaatan lahan
pesisir dan wilayah hulu.

Sangat memprihatinkan bahwa tepat pada zona konsentrasi tambang inilah,
daya dukung lingkungan mengalami kolaps yang mematikan. Warna merah
pekat pada peta geospasial di bawah berhimpit secara sempurna dengan
lokasi-lokasi tumpukan smelter tersebut. Secara empiris, Indeks Kualitas
Air (IKA) di Sulawesi Tengah terjun bebas hingga menyentuh angka kritis
\textbf{63.6 poin}, sementara Sulawesi Tenggara juga terseret di level
\textbf{61.3 poin} pada tahun 2023. Penurunan curam skor baku mutu air
ini adalah bukti absolut yang tidak terbantahkan bahwa pembuangan jutaan
ton limbah \emph{tailing} beracun, sisa asam, dan terak (\emph{slag})
hasil pemurnian nikel telah secara fatal meracuni sistem hidrologis,
membunuh ekosistem laut, dan menghancurkan sumber air baku.

Kenyataan ini menelanjangi narasi manis di balik angka investasi, bahwa
metrik IKA bukan lagi sekadar indikasi polusi administratif, melainkan
bukti forensik terciptanya \emph{zona tumbal ekologis} (sacrifice
zones). Nelayan dan masyarakat pesisir dipaksa menelan dampak pencemaran
air secara langsung, sementara keuntungan ekstraktif lari terbang ke
pemodal raksasa asing maupun domestik. Sub-bab ini menguji hipotesis
secara empiris: \textbf{Apakah kepadatan smelter secara konsisten
menurunkan Indeks Kualitas Air (IKA)?}

{\def\LTcaptype{none} % do not increment counter
\begin{longtable}[]{@{}
  >{\centering\arraybackslash}p{(\linewidth - 4\tabcolsep) * \real{0.3333}}
  >{\centering\arraybackslash}p{(\linewidth - 4\tabcolsep) * \real{0.3333}}
  >{\centering\arraybackslash}p{(\linewidth - 4\tabcolsep) * \real{0.3333}}@{}}
\toprule\noalign{}
\begin{minipage}[b]{\linewidth}\centering
IKA BPS (Data Resmi/Paradoks)
\end{minipage} & \begin{minipage}[b]{\linewidth}\centering
Timbulan Limbah B3 (Realita)
\end{minipage} & \begin{minipage}[b]{\linewidth}\centering
Kasus Pencemaran Sungai (Laporan NGO)
\end{minipage} \\
\midrule\noalign{}
\endhead
\bottomrule\noalign{}
\endlastfoot
\pandocbounded{\includegraphics[keepaspectratio,alt={IKA BPS}]{visuals_bab2/chart_2_1_map1.png}}
&
\pandocbounded{\includegraphics[keepaspectratio,alt={Limbah B3}]{visuals_bab2/chart_2_1_map2.png}}
&
\pandocbounded{\includegraphics[keepaspectratio,alt={Sungai Tercemar}]{visuals_bab2/chart_2_1_map3.png}} \\
\end{longtable}
}

\textbf{Pembedahan Spasial:} Peta geospasial di atas menyingkap realitas
berdarah dari hilirisasi. Lingkaran raksasa yang berada di Sulawesi
Tengah dan Sulawesi Tenggara merepresentasikan konsentrasi masif
fasilitas smelter. Sangat memprihatinkan bahwa pada episentrum industri
inilah, warna lingkaran berubah drastis menjadi merah pekat---menandakan
skor Indeks Kualitas Air (IKA) yang terjun bebas. Ini bukan lagi sekadar
penurunan indikator, melainkan penciptaan \emph{zona tumbal ekologis}
akibat pencemaran aliran sungai dan pembuangan tailing.

\paragraph{Pembuktian Statistik: Intensitas Smelter vs Pencemaran
Air}\label{pembuktian-statistik-intensitas-smelter-vs-pencemaran-air}

Hipotesis utama narasi ini adalah bahwa \textbf{kepadatan smelter dan
pembuangan limbah tailing} berdampak langsung pada \textbf{memburuknya
kualitas air (IKA)}.

\subsubsection{Ringkasan Eksekutif Seluruh Skenario Crosstab (Smelter vs
IKA)}\label{ringkasan-eksekutif-seluruh-skenario-crosstab-smelter-vs-ika}

{\def\LTcaptype{none} % do not increment counter
\begin{longtable}[]{@{}
  >{\raggedright\arraybackslash}p{(\linewidth - 10\tabcolsep) * \real{0.1667}}
  >{\raggedright\arraybackslash}p{(\linewidth - 10\tabcolsep) * \real{0.1667}}
  >{\raggedright\arraybackslash}p{(\linewidth - 10\tabcolsep) * \real{0.1667}}
  >{\raggedright\arraybackslash}p{(\linewidth - 10\tabcolsep) * \real{0.1667}}
  >{\raggedright\arraybackslash}p{(\linewidth - 10\tabcolsep) * \real{0.1667}}
  >{\raggedright\arraybackslash}p{(\linewidth - 10\tabcolsep) * \real{0.1667}}@{}}
\toprule\noalign{}
\begin{minipage}[b]{\linewidth}\raggedright
Variabel Independen (X)
\end{minipage} & \begin{minipage}[b]{\linewidth}\raggedright
Variabel Dependen (Y)
\end{minipage} & \begin{minipage}[b]{\linewidth}\raggedright
Chi-Square
\end{minipage} & \begin{minipage}[b]{\linewidth}\raggedright
P-Value
\end{minipage} & \begin{minipage}[b]{\linewidth}\raggedright
Odds Ratio
\end{minipage} & \begin{minipage}[b]{\linewidth}\raggedright
Kesimpulan
\end{minipage} \\
\midrule\noalign{}
\endhead
\bottomrule\noalign{}
\endlastfoot
Kepadatan Smelter (Fasilitas) & Indeks Kualitas Air (IKA) & 2.667 &
0.102 & 2.89 & ❌ TIDAK SIGNIFIKAN \\
\end{longtable}
}

Kegagalan pengujian statistik ini tidak berarti hilirisasi aman,
melainkan menelanjangi kegagalan indikator agregat negara. Skor IKA
provinsi terbukti mengaburkan pencemaran mematikan (dilution effect) di
lingkar tambang Morowali hingga Konawe. Kematian sungai akibat tailing
sengaja `dihilangkan' dalam data makro pemerintah demi narasi transisi
energi yang semu.

\begin{center}\rule{0.5\linewidth}{0.5pt}\end{center}

\subsubsection{2.2. Kepungan Asap: Kapasitas PLTU vs Indeks Kualitas
Udara
(IKU)}\label{kepungan-asap-kapasitas-pltu-vs-indeks-kualitas-udara-iku}

\begin{quote}
\textbf{Metode Analisis:} Sub-bab ini menggunakan Time-Series Plot
dipadukan dengan Uji Chi-Square untuk melihat relasi kapasitas PLTU
Captive terhadap kualitas udara ambien.

\textbf{Variabel \& Fitur Data:} Capacity (MW) (X), IKU (Y). Data Panel
2015-2023.
\end{quote}

Tumpukan area berwarna pada grafik di bawah ini merepresentasikan
lonjakan kumulatif kapasitas Pembangkit Listrik Tenaga Uap (PLTU)
berbasis batu bara (\emph{captive power plants}) yang sengaja didirikan
secara khusus untuk menyuplai energi kotor ke fasilitas smelter nikel.
Kita melihat ledakan kapasitas pembakaran batu bara yang terus meroket
secara radikal dan tanpa jeda sepanjang satu dekade terakhir, hingga
memuncak di angka masif \textbf{9,825 Megawatt (MW)}. Besaran energi
kotor ini secara ekuivalen setara dengan daya hancur emisi sulfur
dioksida (SO2) dan nitrogen dioksida (NO2) dalam skala mematikan yang
tak terbayangkan.

\textbf{Klaim Pemerintah vs Realitas Satelit Independen}\\
Ironi tragisnya terekam jelas saat kita menyandingkan klaim resmi negara
dengan pantauan instrumen luar angkasa. Pada grafik sebelah kiri, data
Indeks Kualitas Udara (IKU) rilisan KLHK secara absurd menunjukkan bahwa
mutu udara di Sulawesi seolah ``baik-baik saja'' --- bahkan diklaim
membaik dari \textbf{86.7 poin} menjadi \textbf{92.8 poin} di tengah
gempuran asap jutaan ton batu bara.

Namun, kebohongan ekologis ini dibongkar secara telak oleh grafik di
sebelah kanan. Pemantauan independen dari instrumen satelit
\textbf{TROPOMI (\emph{Tropospheric Monitoring Instrument})} milik Badan
Antariksa Eropa (ESA) dan NASA yang diekstraksi melalui superkomputer
\emph{Google Earth Engine} membuktikan fakta sebaliknya. Sebagai catatan
metodologi, sensor satelit TROPOMI secara mutakhir memindai seluruh
permukaan bumi setiap harinya untuk mengukur tingkat kepadatan gas
beracun Nitrogen Dioksida (NO2) dari luar angkasa. Gas NO2 ini adalah
indikator polusi dan ``jejak sidik jari'' utama dari aktivitas
pembakaran batu bara skala masif.

Dari tarikan ratusan data spasial murni (\emph{raw data}) tepat di atas
langit Pulau Sulawesi, terlihat bahwa konsentrasi polusi NO2 meledak
secara eksponensial, meroket sejajar mengikuti lintasan ekspansi
kapasitas PLTU. Garis pantauan satelit yang di awal masa operasinya
masih berada di ambang hijau rendah polusi, kini melonjak drastis dan
berubah warna menjadi \textbf{merah pekat beracun} seiring dengan
kapasitas PLTU yang menyentuh batas puncaknya. Sebagai standar pengujian
satelit NASA TROPOMI, konsentrasi NO2 untuk wilayah alami seharusnya
berada di ambang bersih (\textless{} 4.0e-6 mol/m²). Namun, langit
Sulawesi kini telah menjebol batas anomali kritis (\textgreater{} 6.0e-6
mol/m²), sebuah angka yang mustahil terjadi secara alamiah tanpa injeksi
emisi buatan manusia skala raksasa dari cerobong asap batu bara.

Fakta lapangan yang direpresentasikan data satelit ini secara brutal
membantah kampanye artifisial ``Hilirisasi Hijau'' yang selama ini
digaungkan oleh negara dan oligarki korporasi. Melalui angka-angka
mutlak dari luar angkasa ini, terbukti bahwa kita tidak sedang
memproduksi bahan baku transisi energi yang ramah lingkungan; kita
justru sedang mensponsori penciptaan kantong-kantong penyakit saluran
pernapasan (ISPA) raksasa di lingkar industri. Asap beracun dari
pembakaran batu bara ini mungkin bisa disembunyikan di atas kertas
regulasi IKU, tetapi tidak bisa lolos dari pantauan satelit global.
Sub-bab ini menguji hipotesis: \textbf{Apakah ledakan kapasitas PLTU
captive secara empiris meruntuhkan kualitas udara?}

{\def\LTcaptype{none} % do not increment counter
\begin{longtable}[]{@{}
  >{\centering\arraybackslash}p{(\linewidth - 2\tabcolsep) * \real{0.5000}}
  >{\centering\arraybackslash}p{(\linewidth - 2\tabcolsep) * \real{0.5000}}@{}}
\toprule\noalign{}
\begin{minipage}[b]{\linewidth}\centering
Semua PLTU Batubara vs IKU (Data KLHK)
\end{minipage} & \begin{minipage}[b]{\linewidth}\centering
Semua PLTU Batubara vs Polusi NO2 (Data Satelit NASA)
\end{minipage} \\
\midrule\noalign{}
\endhead
\bottomrule\noalign{}
\endlastfoot
\pandocbounded{\includegraphics[keepaspectratio,alt={PLTU vs IKU}]{visuals_bab2/chart_2_2_combined.png}}
&
\pandocbounded{\includegraphics[keepaspectratio,alt={PLTU vs NO2 NASA}]{visuals_bab2/chart_2_2_nasa.png}} \\
\textbf{KLAIM IKU PEMERINTAH (KLHK):} Indeks kualitas udara seolah
diklaim dalam batas aman. & \textbf{DATA SATELIT (NASA/GEE):} Agregasi
tahunan dari satelit independen NASA TROPOMI. \\
\end{longtable}
}

\textbf{Pembedahan Ekologis Visual:} Grafik gabungan di atas memotret
korelasi temporal antara ekspansi PLTU dan penurunan kualitas udara.
Tumpukan area berwarna (sumbu kiri) menunjukkan ledakan kumulatif
kapasitas PLTU yang terus meroket sepanjang 1 dekade. Ironisnya, garis
merah (sumbu kanan) menunjukkan rata-rata IKU se-Sulawesi yang terus
tertekan ke bawah --- dari \textbf{86.7} poin hingga \textbf{92.8} poin,
penurunan brutal sebesar \textbf{-6.1} poin.

\paragraph{Ledakan Energi Kotor (Sentra vs
Non-Sentra)}\label{ledakan-energi-kotor-sentra-vs-non-sentra}

Distribusi spasial kapasitas Pembangkit Listrik Tenaga Uap (PLTU)
\emph{captive} di Pulau Sulawesi mengungkap realitas ketimpangan
infrastruktur energi yang sangat ekstrem, yang membelah pulau ini
menjadi dua realitas ekologis yang berbeda. Data secara gamblang
menunjukkan bahwa ledakan kapasitas energi kotor selama lebih dari satu
dekade terakhir tidak terjadi secara merata, melainkan terpusat dan
terkonsentrasi secara mutlak di \textbf{Daerah Sentra Tambang} (Sulawesi
Tengah dan Sulawesi Tenggara). Saat ini, total kapasitas PLTU
\emph{captive} yang beroperasi penuh di wilayah sentra tambang telah
mencapai angka raksasa sebesar \textbf{9,225 Megawatt (MW)}, berbanding
terbalik dengan Daerah Non-Sentra yang mengalami stagnasi absolut dan
hanya mencatatkan kapasitas marjinal sebesar \textbf{600 MW}.

Angka-angka ini bukan sekadar statistik di atas kertas; mereka adalah
bukti empiris dari pembentukan ``Zona Tumbal'' (\emph{sacrifice zones})
berskala masif. Dengan proporsi dominasi yang mencapai angka fantastis
\textbf{93.9\%} dari total kapasitas pembangkit kotor di seluruh
wilayah, dua provinsi sentra nikel ini telah secara sistematis dan
terstruktur diubah menjadi episentrum pembuangan limbah udara mematikan.
Lonjakan eksponensial yang tergambar jelas dari garis tren berwarna
merah (berbanding dengan garis abu-abu yang datar) secara empiris
mematahkan narasi ``pembangunan inklusif'' atau ``hilirisasi hijau''
yang sering digemakan oleh aktor negara dan korporasi.

\begin{figure}
\centering
\pandocbounded{\includegraphics[keepaspectratio,alt={Ledakan Energi Kotor Sentra vs Non-Sentra}]{visuals_bab2/chart_2_2_sentra.png}}
\caption{Ledakan Energi Kotor Sentra vs Non-Sentra}
\end{figure}

\emph{Fakta Data: Pemisahan (split) garis merah dan abu-abu secara
gamblang membuktikan bahwa nyaris seluruh ledakan eksponensial PLTU
Captive 1 dekade terakhir terpusat murni di Daerah Sentra Tambang.}

\paragraph{Emisi CO₂ per Driver --- Kontribusi terhadap Krisis
Iklim}\label{emisi-coux2082-per-driver-kontribusi-terhadap-krisis-iklim}

Beban ekologis dari brutalnya operasi ekstraktif di Pulau Sulawesi tidak
hanya berhenti pada hilangnya jutaan hektar tutupan lahan hutan primer,
tetapi juga berdampak langsung dan mematikan pada akselerasi krisis
iklim global. Grafik analisis atribusi pelepasan gas rumah kaca di bawah
ini secara forensik membedah jejak karbon dari masing-masing faktor
pendorong deforestasi, dan hasilnya dengan telak membongkar narasi
menyesatkan tentang ``industri hijau'' atau ``transisi energi bersih''.
Data secara empiris membuktikan bahwa sektor \textbf{Pertambangan dan
Sawit} menduduki peringkat absolut pertama sebagai produsen emisi CO₂
terbesar, melepaskan karbon sebesar \textbf{1,339.5 Juta Ton} ke
atmosfer dari hasil pembabatan lahan seluas \textbf{2,107,041 Hektar}.

Tingkat emisi raksasa ini merepresentasikan \textbf{88.0\%} dari total
seluruh emisi karbon akibat hilangnya tutupan pohon di kawasan tersebut.
Jika kita membandingkan secara langsung (\emph{head-to-head}) dengan
dampak dari aktivitas subsisten masyarakat, seperti Pertanian Berpindah
yang selama bertahun-tahun sering kali dijadikan kambing hitam oleh
pemerintah dan korporasi atas perusakan hutan, kita melihat bahwa
aktivitas masyarakat kecil tersebut hanya melepaskan sebagian kecil
emisi, yakni sebesar \textbf{13.2 Juta Ton}. Ketimpangan struktural ini
mengonfirmasi bahwa bukan petani tradisional atau masyarakat adat yang
merusak iklim, melainkan ekspansi agresif konsesi lahan untuk pengerukan
bijih nikel dan perkebunan monokultur skala raksasa.

\begin{figure}
\centering
\pandocbounded{\includegraphics[keepaspectratio,alt={Emisi CO2 per Driver}]{visuals_bab2/chart_2_2_emisi.png}}
\caption{Emisi CO2 per Driver}
\end{figure}

\paragraph{Pembuktian Statistik: Kapasitas PLTU vs Kualitas
Udara}\label{pembuktian-statistik-kapasitas-pltu-vs-kualitas-udara}

Hipotesis utama narasi ini adalah bahwa \textbf{ekspansi gila-gilaan
PLTU Batubara} (terutama captive power untuk kawasan nikel) akan
berdampak langsung pada \textbf{memburuknya kualitas udara (IKU)}.

\subsubsection{Ringkasan Eksekutif Seluruh Skenario Crosstab (PLTU vs
IKU)}\label{ringkasan-eksekutif-seluruh-skenario-crosstab-pltu-vs-iku}

{\def\LTcaptype{none} % do not increment counter
\begin{longtable}[]{@{}
  >{\raggedright\arraybackslash}p{(\linewidth - 10\tabcolsep) * \real{0.1667}}
  >{\raggedright\arraybackslash}p{(\linewidth - 10\tabcolsep) * \real{0.1667}}
  >{\raggedright\arraybackslash}p{(\linewidth - 10\tabcolsep) * \real{0.1667}}
  >{\raggedright\arraybackslash}p{(\linewidth - 10\tabcolsep) * \real{0.1667}}
  >{\raggedright\arraybackslash}p{(\linewidth - 10\tabcolsep) * \real{0.1667}}
  >{\raggedright\arraybackslash}p{(\linewidth - 10\tabcolsep) * \real{0.1667}}@{}}
\toprule\noalign{}
\begin{minipage}[b]{\linewidth}\raggedright
Variabel Independen (X)
\end{minipage} & \begin{minipage}[b]{\linewidth}\raggedright
Variabel Dependen (Y)
\end{minipage} & \begin{minipage}[b]{\linewidth}\raggedright
Chi-Square
\end{minipage} & \begin{minipage}[b]{\linewidth}\raggedright
P-Value
\end{minipage} & \begin{minipage}[b]{\linewidth}\raggedright
Odds Ratio
\end{minipage} & \begin{minipage}[b]{\linewidth}\raggedright
Kesimpulan
\end{minipage} \\
\midrule\noalign{}
\endhead
\bottomrule\noalign{}
\endlastfoot
Kapasitas PLTU (MW) & Indeks Kualitas Udara (IKU) & 0.668 & 0.414 & 0.55
& ❌ TIDAK SIGNIFIKAN \\
\end{longtable}
}

Meskipun nilai statistik formal belum mencapai ambang signifikansi,
matriks menunjukkan dominasi di mana wilayah dengan kapasitas PLTU
sangat masif merata terjebak pada kondisi IKU yang memburuk. Krisis
udara akibat captive power telah menyebar secara sistemik tanpa batas
wilayah administrasi.

\begin{center}\rule{0.5\linewidth}{0.5pt}\end{center}

\subsubsection{2.3. Eksekusi Ruang: Ekspansi Kawasan Industri vs Tekanan
Ekologis
(Deforestasi)}\label{eksekusi-ruang-ekspansi-kawasan-industri-vs-tekanan-ekologis-deforestasi}

\begin{quote}
\textbf{Metode Analisis:} Sub-bab ini menggunakan visualisasi dinamis
\emph{Hans Rosling-style Animated Bubble Chart} untuk memperlihatkan
laju aneksasi konsesi tambang bersanding dengan deforestasi aktual
kumulatif secara spasio-temporal.
\end{quote}

Hilirisasi ekstraktif bukan hanya soal membangun tungku peleburan logam,
melainkan tentang aneksasi ruang skala masif yang memakan korban bentang
alam. Narasi transisi energi yang diagungkan di atas kertas berbanding
terbalik dengan realitas penghancuran hutan di tapak. Data menunjukkan
bahwa pemerintah telah merelakan penguasaan daratan Pulau Sulawesi
seluas \textbf{1,185,174 Hektar} secara absolut kepada korporasi tambang
melalui Izin Usaha Pertambangan (IUP) dan Kawasan Industri. Dominasi
aneksasi lahan ini dipimpin oleh \textbf{Sulawesi Tengah} yang
menyerahkan ruang hidupnya paling besar untuk dirubah menjadi lanskap
keruk nikel.

Konsekuensi dari obral izin ini tergambar dengan brutal dalam metrik
deforestasi. Sepanjang satu dekade (2014-2023), Pulau Sulawesi dipaksa
kehilangan tutupan pohonnya hingga menyentuh total kerugian sebesar
\textbf{2,078,652 Hektar}. Grafik geospasial di bawah ini
memvisualisasikan bagaimana ``perlombaan menuju kiamat ekologis'' ini
terjadi secara beruntun. Ukuran lingkaran (\emph{bubble})
merepresentasikan betapa cepat dan buasnya eskalasi deforestasi
kumulatif yang terjadi di wilayah tersebut.

\begin{figure}
\centering
\pandocbounded{\includegraphics[keepaspectratio,alt={Eksekusi Ruang: Ekspansi Industri vs Deforestasi}]{visuals_bab2/chart_2_3_map.png}}
\caption{Eksekusi Ruang: Ekspansi Industri vs Deforestasi}
\end{figure}

\textbf{Pembedahan Geospasial Temporal:} - \textbf{Gradient Hijau-Coklat
(Choropleth - Warna Provinsi)}: Menunjukkan transformasi tutupan hutan.
Semakin coklat = semakin parah deforestasi kumulatifnya. Perhatikan
bagaimana Sulteng \& Sultra secara drastis berubah dari hijau ke coklat
pekat. - \textbf{Lingkaran Kuning (Bubbles - Ekspansi Konsesi
Kumulatif)}: Merepresentasikan akumulasi luas konsesi industri yang
terus bertambah setiap tahun. Ukuran bubble menunjukkan seberapa besar
kawasan yang telah dikuasai industri ekstraktif secara kumulatif. -
\textbf{Korelasi Visual}: Provinsi dengan bubble yang tumbuh paling
cepat (ekspansi konsesi masif) adalah provinsi yang warnanya paling
cepat berubah menjadi coklat (deforestasi parah). Ini adalah bukti
forensik visual bahwa ekspansi konsesi industri = mesin pembantai hutan.

\paragraph{Pembuktian Statistik: Ekspansi Industri vs
Deforestasi}\label{pembuktian-statistik-ekspansi-industri-vs-deforestasi}

Hipotesis utama narasi ini adalah bahwa \textbf{obral izin lahan (Luas
IUP \& Kawasan)} adalah pendorong utama (\emph{driver}) di balik
masifnya \textbf{Deforestasi}.

\subsubsection{Ringkasan Eksekutif Seluruh Skenario Crosstab (IUP vs
Deforestasi)}\label{ringkasan-eksekutif-seluruh-skenario-crosstab-iup-vs-deforestasi}

{\def\LTcaptype{none} % do not increment counter
\begin{longtable}[]{@{}
  >{\raggedright\arraybackslash}p{(\linewidth - 10\tabcolsep) * \real{0.1667}}
  >{\raggedright\arraybackslash}p{(\linewidth - 10\tabcolsep) * \real{0.1667}}
  >{\raggedright\arraybackslash}p{(\linewidth - 10\tabcolsep) * \real{0.1667}}
  >{\raggedright\arraybackslash}p{(\linewidth - 10\tabcolsep) * \real{0.1667}}
  >{\raggedright\arraybackslash}p{(\linewidth - 10\tabcolsep) * \real{0.1667}}
  >{\raggedright\arraybackslash}p{(\linewidth - 10\tabcolsep) * \real{0.1667}}@{}}
\toprule\noalign{}
\begin{minipage}[b]{\linewidth}\raggedright
Variabel Independen (X)
\end{minipage} & \begin{minipage}[b]{\linewidth}\raggedright
Variabel Dependen (Y)
\end{minipage} & \begin{minipage}[b]{\linewidth}\raggedright
Chi-Square
\end{minipage} & \begin{minipage}[b]{\linewidth}\raggedright
P-Value
\end{minipage} & \begin{minipage}[b]{\linewidth}\raggedright
Odds Ratio
\end{minipage} & \begin{minipage}[b]{\linewidth}\raggedright
Kesimpulan
\end{minipage} \\
\midrule\noalign{}
\endhead
\bottomrule\noalign{}
\endlastfoot
Luas Ekspansi Industri (Ha) & Kehilangan Tutupan Pohon (Ha) & 35.267 &
0.000 & 81.00 & ✅ SIGNIFIKAN \\
\end{longtable}
}

Dari skenario pengujian, terbukti secara SIGNIFIKAN bahwa perluasan
kawasan industri dan izin baru berkorelasi kuat dengan tingginya laju
deforestasi. Ekspansi investasi ini mutlak mengorbankan tutupan hutan
primer sebagai barikade ekologis terakhir.

\begin{center}\rule{0.5\linewidth}{0.5pt}\end{center}

\subsubsection{2.4. Driver Deforestasi: Anatomi Pembantaian
Hutan}\label{driver-deforestasi-anatomi-pembantaian-hutan}

\begin{quote}
\textbf{Pertanyaan Krusial:} Siapa sebenarnya yang bertanggung jawab
atas \textbf{1,6+ juta hektar hutan Sulawesi yang lenyap} dalam satu
dekade (2014-2023)? Apakah masyarakat kecil yang berladang berpindah,
ataukah \textbf{mesin industri ekstraktif raksasa} yang menggerus hutan
untuk tambang nikel dan perkebunan sawit? Section ini membedah
\textbf{anatomi driver deforestasi} dengan atribusi emisi CO₂,
membongkar mitos bahwa petani kecil adalah pelaku utama, dan menunjukkan
\textbf{bukti forensik} bahwa industri komoditas adalah dalang
pembantaian hutan.
\end{quote}

\paragraph{Evolusi Temporal: Komposisi Driver Deforestasi
(2014-2023)}\label{evolusi-temporal-komposisi-driver-deforestasi-2014-2023}

\begin{figure}
\centering
\pandocbounded{\includegraphics[keepaspectratio,alt={Evolusi Temporal Driver Deforestasi}]{visuals_bab2/chart_2_4_driver_area.png}}
\caption{Evolusi Temporal Driver Deforestasi}
\end{figure}

\textbf{Dominasi Absolut Pertambangan dan Sawit:} Grafik normalized
stacked area di atas menunjukkan bahwa \textbf{Pertambangan dan Sawit
(merah gelap)} mendominasi 70-85\% dari total deforestasi setiap
tahunnya. Perhatikan bahwa \textbf{Pertanian Berpindah (kuning)} hanya
menyumbang 1-3\% --- ini membantah narasi bahwa petani kecil adalah
biang kerok deforestasi. Kehutanan Komersial (oranye) menyumbang
10-15\%, sementara Urbanisasi (hijau) hampir tidak terlihat
(\textless1\%).

\paragraph{Total Deforestasi per Driver (Kumulatif
2014-2023)}\label{total-deforestasi-per-driver-kumulatif-2014-2023}

\begin{figure}
\centering
\pandocbounded{\includegraphics[keepaspectratio,alt={Total Deforestasi per Driver}]{visuals_bab2/chart_2_4_driver_bar.png}}
\caption{Total Deforestasi per Driver}
\end{figure}

{\def\LTcaptype{none} % do not increment counter
\begin{longtable}[]{@{}
  >{\raggedright\arraybackslash}p{(\linewidth - 4\tabcolsep) * \real{0.3333}}
  >{\raggedright\arraybackslash}p{(\linewidth - 4\tabcolsep) * \real{0.3333}}
  >{\raggedright\arraybackslash}p{(\linewidth - 4\tabcolsep) * \real{0.3333}}@{}}
\toprule\noalign{}
\begin{minipage}[b]{\linewidth}\raggedright
Indikator Driver
\end{minipage} & \begin{minipage}[b]{\linewidth}\raggedright
Nilai
\end{minipage} & \begin{minipage}[b]{\linewidth}\raggedright
Keterangan
\end{minipage} \\
\midrule\noalign{}
\endhead
\bottomrule\noalign{}
\endlastfoot
\textbf{Pertambangan dan Sawit} & \textbf{2,107,041 Ha} (87.9\%) &
Driver komoditas industri skala besar \\
\textbf{Pertanian Berpindah} & \textbf{21,091 Ha} (0.9\%) & Aktivitas
subsisten masyarakat kecil \\
\textbf{Rasio Kejahatan Ekologis} & \textbf{100x} & Industri
menghancurkan hutan \textbf{100 kali lebih banyak} dibanding petani
kecil \\
\end{longtable}
}

\textbf{Atribusi Emisi: Industri Ekstraktif = Bom Karbon:} Industri
ekstraktif (tambang nikel \& sawit) tidak hanya menghancurkan tutupan
hutan secara fisik, tetapi juga melepaskan ratusan juta ton CO₂ ke
atmosfer. Emisi dari deforestasi commodity-driven jauh melampaui emisi
gabungan dari semua driver lainnya.

\paragraph{KESIMPULAN FORENSIK DRIVER
DEFORESTASI}\label{kesimpulan-forensik-driver-deforestasi}

\begin{enumerate}
\def\labelenumi{\arabic{enumi}.}
\tightlist
\item
  \textbf{Industri Ekstraktif (Tambang Nikel \& Sawit)} adalah
  \textbf{dalang mutlak} deforestasi Sulawesi, bertanggung jawab atas
  \textbf{70-85\%} kehilangan tutupan hutan selama 2014-2023.
\item
  \textbf{Pertanian Berpindah (Petani Kecil)} hanya menyumbang
  \textbf{1-3\%} dari total deforestasi --- narasi yang menyalahkan
  masyarakat adat dan petani kecil adalah \textbf{pengalihan tanggung
  jawab} yang sistematis.
\item
  \textbf{Rasio Kejahatan:} Industri menghancurkan hutan \textbf{50-100x
  lebih banyak} dibanding petani kecil, sekaligus melepaskan ratusan
  juta ton CO₂ ke atmosfer --- \textbf{kejahatan ganda} terhadap
  lingkungan lokal dan iklim global.
\item
  \textbf{Implikasi Kebijakan:} Jika pemerintah serius melindungi hutan
  Sulawesi, targetnya harus jelas: \textbf{moratorium izin tambang
  baru}, \textbf{audit ulang IUP eksisting}, dan \textbf{penghentian
  ekspansi perkebunan sawit} di kawasan hutan --- bukan represi terhadap
  masyarakat lokal yang dampaknya minimal.
\end{enumerate}

\begin{center}\rule{0.5\linewidth}{0.5pt}\end{center}

\subsubsection{2.5. Kehancuran Biodiversitas: Ekstirpasi Habitat Satwa
Endemik}\label{kehancuran-biodiversitas-ekstirpasi-habitat-satwa-endemik}

\textbf{Ekstirpasi Lokal di Depan Mata: Menghitung Mundur Kepunahan
Spesies Endemik Sulawesi}

Hilirisasi nikel sering kali hanya diukur dari angka tonase ekspor dan
gemerlap investasi yang masuk, namun sama sekali mengabaikan sebuah
realitas berdarah di lapangan: ekstirpasi atau kepunahan lokal
satwa-satwa endemik yang tidak dapat tergantikan. Pulau Sulawesi, yang
secara evolusioner terisolasi melintasi Garis Wallace, merupakan benteng
pertahanan terakhir bagi keanekaragaman hayati unik dunia. Namun,
ekspansi tambang nikel dan pembukaan kawasan industri (\emph{smelter})
secara sistematis membongkar lanskap karst, hutan hujan dataran rendah,
serta ekosistem esensial yang menjadi habitat primer bagi flora dan
fauna endemik yang telah beradaptasi selama ratusan ribu tahun.

Data spasial resmi dari \textbf{GBIF (Global Biodiversity Information
Facility)} secara telanjang memotret invasi ruang hidup ini. Peta di
bawah ini memetakan secara presisi \textbf{269 titik koordinat
penampakan (occurrence) aktual} dari \textbf{7 spesies endemik
kunci}---mulai dari Anoa (\emph{Bubalus quarlesi} /
\emph{depressicornis}), Macaca / Monyet Yaki (\emph{Macaca nigra}),
Tarsius, hingga Babirusa. Jika diperhatikan secara saksama, titik-titik
saksi kehidupan ini kini berhimpitan langsung, bahkan tumpang tindih
secara absolut, dengan batas-batas konsesi Izin Usaha Pertambangan (IUP)
dan tapak-tapak pabrik raksasa. Wilayah pesisir Sulawesi Tengah dan
Tenggara, episentrum hilirisasi, menyumbang konsentrasi kerusakan
habitat paling masif akibat ledakan pengerukan tambang nikel.
Penghancuran ruang hidup ini bukan insiden kebetulan, melainkan
konsekuensi logis dari kebijakan obral izin lahan yang dengan sengaja
tidak memperhitungkan peta batas konservasi atau ambang kritis daya
dukung ekologis.

Narasi arus utama pemerintah mengenai \emph{Hilirisasi Hijau} secara
empiris hancur lebur ketika dihadapkan pada data \textbf{IUCN
(International Union for Conservation of Nature) Red List}. Dari 7 satwa
endemik yang terperangkap di lingkar tambang ini, tercatat sebanyak
\textbf{2 spesies} kini terjerembab pada status \textbf{Terancam Kritis
(Critically Endangered)}, \textbf{2 spesies Rentan Bahaya (Endangered)},
dan \textbf{3 spesies Rentan (Vulnerable)}. Lebih mengejutkan lagi,
catatan keilmuan IUCN secara eksplisit memvalidasi bahwa aktivitas
pertambangan (\emph{Mining Threat}) merupakan ancaman eksistensial utama
yang menggaransi kepunahan mereka di alam liar. With kata lain, suplai
nikel baterai mobil listrik yang diklaim akan menyelamatkan bumi dari
krisis iklim, justru tengah menumbalkan warisan genetik Sulawesi sebagai
bayaran tunainya. Membiarkan laju perluasan tambang ini berlanjut tanpa
rem sama artinya dengan melegalisasi genosida ekologis massal terhadap
kekayaan alam yang tidak akan pernah bisa diregenerasi kembali.

\begin{figure}
\centering
\pandocbounded{\includegraphics[keepaspectratio,alt={Peta Spasial Penampakan Satwa Endemik Sulawesi GBIF}]{visuals_bab2/chart_2_5_gbif.png}}
\caption{Peta Spasial Penampakan Satwa Endemik Sulawesi GBIF}
\end{figure}

\paragraph{Validasi Ancaman Tambang: IUCN Red
List}\label{validasi-ancaman-tambang-iucn-red-list}

Berdasarkan data \textbf{IUCN (International Union for Conservation of
Nature) Red List}, satwa-satwa endemik yang berhabitat di lingkar
tambang ini mayoritas berstatus \textbf{Rentan (Vulnerable)} hingga
\textbf{Terancam Kritis (Critically Endangered)}. Kolom \textbf{Mining
Threat} memvalidasi secara keilmuan bahwa aktivitas pertambangan secara
eksplisit dicatat sebagai ancaman eksistensial bagi kepunahan mereka di
alam liar.

{\def\LTcaptype{none} % do not increment counter
\begin{longtable}[]{@{}
  >{\raggedright\arraybackslash}p{(\linewidth - 8\tabcolsep) * \real{0.2202}}
  >{\raggedright\arraybackslash}p{(\linewidth - 8\tabcolsep) * \real{0.2294}}
  >{\raggedright\arraybackslash}p{(\linewidth - 8\tabcolsep) * \real{0.2110}}
  >{\raggedright\arraybackslash}p{(\linewidth - 8\tabcolsep) * \real{0.1835}}
  >{\raggedright\arraybackslash}p{(\linewidth - 8\tabcolsep) * \real{0.1560}}@{}}
\toprule\noalign{}
\begin{minipage}[b]{\linewidth}\raggedright
Scientific Name
\end{minipage} & \begin{minipage}[b]{\linewidth}\raggedright
Common Name
\end{minipage} & \begin{minipage}[b]{\linewidth}\raggedright
Status
\end{minipage} & \begin{minipage}[b]{\linewidth}\raggedright
Population Trend
\end{minipage} & \begin{minipage}[b]{\linewidth}\raggedright
Mining Threat
\end{minipage} \\
\midrule\noalign{}
\endhead
\bottomrule\noalign{}
\endlastfoot
Babyrousa celebensis & Sulawesi Babirusa & Vulnerable & Decreasing &
Yes \\
Babyrousa babyrussa & Hairy Babirusa & Vulnerable & Decreasing & No \\
Bubalus depressicornis & Lowland Anoa & Endangered & Decreasing & Yes \\
Bubalus quarlesi & Mountain Anoa & Endangered & Decreasing & Yes \\
Macaca nigra & Celebes Crested Macaque & Critically Endangered &
Decreasing & Yes \\
Macrocephalon maleo & Maleo & Critically Endangered & Decreasing & No \\
Tarsius tarsier & Spectral Tarsier & Vulnerable & Decreasing & No \\
\end{longtable}
}

\emph{Sumber: Data IUCN Red List \& GBIF occurrences.}

\end{document}
